% Generated draft: requires mathematical and source review.
\documentclass[11pt]{article}
\usepackage[margin=25mm]{geometry}
\usepackage{fontspec}
\setmainfont{Latin Modern Roman}
\usepackage{amsmath,amssymb,mathtools,unicode-math}
\setmathfont{Latin Modern Math}
\usepackage{enumitem,needspace,xurl,hyperref}
\setlength{\parindent}{0pt}
\setlength{\parskip}{5pt}
\setlength{\emergencystretch}{3em}
\newcommand{\N}{\mathbb{N}}
\newcommand{\Z}{\mathbb{Z}}
\newcommand{\Q}{\mathbb{Q}}
\newcommand{\R}{\mathbb{R}}
\newcommand{\C}{\mathbb{C}}
\newcommand{\F}{\mathbb{F}}
\newcommand{\A}{\mathbb{A}}
\newcommand{\PP}{\mathbb{P}}
\newcommand{\E}{\mathbb{E}}
\newcommand{\eps}{\varepsilon}
\newcommand{\dd}{\,\mathrm{d}}
\begin{document}
\section*{UnsolvedMath ID: 10400050}
\textbf{Source problem code:} catalogue entry
\subsection{Definitions and mathematical statement}
Problem 2.28, attributed to M.-J. Jeong and C.-Y. Park, asks for a minimal finite subset \(A_n\) of \(V_n\) satisfying
\[
 A_n\subset V_n,\qquad |A_n|<\infty,\qquad
 \operatorname{span}(A_n)=V_n.
\]
Neither \(V_n\), the coefficient system for the span, nor the range of \(n\) is defined in the supplied record. Thus the object whose generators are requested cannot be identified precisely from this extract.

If \(V_n\) is intended as a vector space over a field \(F\), then
\[
 \operatorname{span}_F(A_n)=
 \left\{\sum_{j=1}^{r}a_jv_j:r\ge0,\ a_j\in F,\ v_j\in A_n\right\}.
\]
In that setting an inclusion-minimal spanning set satisfies \(\operatorname{span}_F(B)\ne V_n\) for every proper subset \(B\subsetneq A_n\); equivalently it is a basis. A minimum-cardinality spanning set would have size \(\dim_F V_n\). These conditional interpretations explain ``minimal'' but do not identify \(V_n\) or provide the requested set. If the source instead uses a module or a constrained family of invariants, its spanning and minimality conventions must be stated explicitly. No particular finite-type invariant space, coefficient field, or uniform construction in \(n\) is inferred from the notation alone.
\subsection{Short English statement}
Find a smallest or irredundant finite spanning family for the source's space \(V_n\). Its definition and coefficients are missing, so a source-specific generating family cannot yet be specified. The supplied open-triage summary calls the question not obviously resolved.
\subsection{Sources}
\begin{itemize}
\item Local UnsolvedMath registry, permanent ID \(10400050\), exact source code \texttt{catalogue entry}, Problem 2.28, attributed in the statement to M.-J. Jeong and C.-Y. Park. The supplied title is truncated after ``(M.-J.''
\item Ohtsuki, \emph{Problems on invariants of knots and 3-manifolds} (2002), Problem 2.28, PDF page 51: \url{https://msp.org/gtm/2002/04/gtm-2002-04-024s.pdf}. Bold printed headings distinguish 231 numbered items from references in surrounding exposition.
\item Extraction pdf-font-structure, accessed 2026-07-29. Source evidence NEEDS\_REVIEW presents an open problem and requires a current primary-literature audit. Research classification OPEN-TRIAGE; supplied summary: minimal generating set for \(V_n\), not obviously resolved. No further authoritative reference defining \(V_n\) is supplied.
\item Dataset attribution: Ulam AI / UnsolvedMath Contributors, CC BY 4.0. Source rights NEEDS\_REVIEW: publicly accessible author/publisher copy, redistribution terms requiring review.
\end{itemize}
\end{document}
